\chapter{Fejlesztői dokumentáció}



\section{Rendszerkövetelmények}

\subsection{Fejlesztői környezet}

\begin{itemize}
    \item \textbf{Python 3.8+} - Adatfeldolgozás, modellezés és háttérlogika
    \item \textbf{Node.js 16.x+} - Frontend fejlesztés
    \item \textbf{VSCode} vagy más fejlesztőkörnyezet
\end{itemize}

\subsection{Függőségek}

\begin{itemize}
    \item \textbf{Python csomagok:}
    \begin{itemize}
        \item pandas>=1.2.0
        \item numpy>=1.20.0
        \item scikit-learn>=0.24.0
        \item torch>=1.8.0
        \item xgboost>=1.3.0
        \item matplotlib>=3.3.0
        \item seaborn>=0.11.0
        \item requests>=2.25.0
    \end{itemize}
    
    \item \textbf{Node csomagok:}
    \begin{itemize}
        \item next>=12.0.0
        \item react>=17.0.0
        \item tailwindcss>=3.0.0
        \item typescript>=4.5.0
    \end{itemize}
\end{itemize}

\section{Architekturális áttekintés}

A Premier League Predictor modularizált, többrétegű architektúrát használ, amely elválasztja az adatgyűjtési, adatfeldolgozási, modellezési és felhasználói felület komponenseket. Az alábbi ábra szemlélteti a rendszer fő komponenseit és kapcsolataikat:

\begin{figure}[H]
    \centering    	\includegraphics[width=1\textwidth, frame]{images/systemarchitecture.png}
    \caption{A rendszer architektúrája}
    \label{fig:system_architecture}
\end{figure}

\textbf{Fő komponensek:}

\begin{enumerate}
    \item \textbf{Adatgyűjtés (scraping modul)} - FBref.com oldalról történő adatgyűjtés, nyers adatok mentése
    \item \textbf{Adatelőkészítés (data\_prep modul)} - Nyers adatok átalakítása, jellemzők kinyerése és tanító/teszt adatok létrehozása
    \item \textbf{Modellezés (train modul)} - Különböző predikciós modellek implementálása és tanítása
    \item \textbf{Előrejelzés (predict modul)} - Tanított modellek használata új mérkőzések kimenetelének előrejelzésére
    \item \textbf{Modell összehasonlítás (compare modul)} - Modellek teljesítményének összehasonlítása
    \item \textbf{Felhasználói felület (football-predictor-ui)} - Next.js alapú webes felület
\end{enumerate}

\section{Adatmodellek és adatszerkezetek}

\subsection{Adatfolyam áttekintés}

A rendszer adatfeldolgozási folyamata az alábbi fő lépésekből áll:

\begin{enumerate}
    \item \textbf{Nyers adatok begyűjtése} - A \texttt{scraping.py} segítségével történik adatgyűjtés a FBref.com oldalról
    \item \textbf{Team Performance dataset létrehozása} - A \texttt{create\_team\_performance\_dataset.py} egyesíti a különböző statisztikákat
    \item \textbf{Modell adatok előkészítése} - A \texttt{prepare\_model\_data.py} párosítja a csapatstatisztikákat a mérkőzésadatokkal
    \item \textbf{Adatok felosztása} - Időrendi alapon történik a tanító-teszt adatok felosztása (a 2023-2024 és 2024-2025 szezonok teszt adatként szerepelnek)
    \item \textbf{Modell-specifikus adatelőkészítés} - Minden modelltípus sajátos adatelőkészítést igényel (pl. normalizálás, kódolás)
\end{enumerate}

\subsection{Nyers adatok}

A rendszer az alábbi nyers adatokat használja, amelyek a \texttt{data/raw} mappában találhatók:

\begin{itemize}
    \item \texttt{all\_seasons\_standard\_stats.csv}
    \item \texttt{all\_seasons\_passing.csv}
    \item \texttt{all\_seasons\_shooting.csv}
    \item \texttt{all\_seasons\_defensive.csv}
    \item \texttt{all\_seasons\_possession.csv}
    \item \texttt{all\_seasons\_goal\_creation.csv} 
    \item \texttt{all\_seasons\_goalkeeping.csv} 
    \item \texttt{all\_seasons\_league\_table.csv} 
    \item \texttt{premier\_league\_match\_outcomes\_2017-2018\_to\_2024-2025.csv}
\end{itemize}


\subsection{Feldolgozott adatok}

Az előkészített adatok a \texttt{data/processed} és \texttt{data/train\_test} mappákban találhatók:

\begin{itemize}
    \item \texttt{team\_performance\_dataset.csv} - Egyesített csapatteljesítmény-adatok
    \item \texttt{train\_data\_chronological.csv} - Időrendben felosztott tanítási adatkészlet
    \item \texttt{test\_data\_chronological.csv} - Időrendben felosztott tesztelési adatkészlet
\end{itemize}

A \texttt{team\_performance\_dataset.csv} fájl tartalmazza az egyes csapatok szezonszintű teljesítményét, beleértve:
\begin{itemize}
    \item Alapvető teljesítménymutatók (győzelmek, döntetlenek, vereségek)
    \item Hazai és vendég teljesítmény külön-külön
    \item Lőtt és kapott gólok
    \item Lövési, passzolási és birtoklási statisztikák
    \item Derivált mutatók, mint például a győzelmi arány, átlagos szerzett gólok száma, stb.
\end{itemize}

A modell betanításához használt adatkészletek (\texttt{train\_data\_chronological.csv} és \texttt{test\_data\_chronological.csv}) párosítják a csapatstatisztikákat a konkrét mérkőzésekkel, és a következő típusú jellemzőket tartalmazzák:
\begin{itemize}
    \item \texttt{Home\_*} - Hazai csapat statisztikái (pl. Home\_MP, Home\_W, Home\_GF, stb.)
    \item \texttt{Away\_*} - Vendég csapat statisztikái (pl. Away\_MP, Away\_W, Away\_GF, stb.)
    \item \texttt{Diff\_*} - A hazai és vendég csapat statisztikái közötti különbségek
\end{itemize}

A célváltozó minden esetben a \texttt{Result} mező, amely lehet "H" (hazai győzelem), "D" (döntetlen) vagy "A" (vendég győzelem).

\section{Modulok részletes leírása}

\subsection{Scraping modul}

A \texttt{scraping.py} script felelős az adatgyűjtésért a FBref.com oldalról. Két fő funkciót tartalmaz:

\begin{itemize}
    \item \texttt{scrape\_detailed\_stats(season)} - Egy adott szezon részletes statisztikáinak gyűjtése
    \item \texttt{scrape\_multiple\_seasons(seasons)} - Több szezon adatainak összegyűjtése és kombinálása
    \item \texttt{scrape\_match\_outcomes(seasons)} - Mérkőzéseredmények gyűjtése
\end{itemize}

A script a ScraperAPI-t használja a lekérések kezelésére az oldal rate limiting elkerülése érdekében. Ennek a használatához saját ScraperAPI tokenre van szükség, amit a "\texttt{*api\_key*}" helyekre kell beilleszteni a kódban.
Amennyiben egy új szezon adatait is szeretnénk összegyűjteni az oldalról, a \texttt{\_\_main\_\_} metódusban a \texttt{seasons\_to\_scrape} listához lehet hozzáírni \textit{"2024-2025"} formátumban.
Az összegyűjtött adatokat CSV fájlokba menti.

Futtatás:
\begin{lstlisting}[language=bash]
    python scraping/scraping.py
\end{lstlisting}

\subsection{Data Prep modul}

Az adatelőkészítést két fő script végzi:

\begin{itemize}
    

\item\texttt{create\_team\_performance\_dataset.py} - A különböző statisztikai adatokat tartalmazó CSV fájlokat egyesíti egy \texttt{team\_performance.csv} fájlba, amely tartalmazza az összes releváns jellemzőt minden csapathoz minden szezonban.

\item\texttt{prepare\_model\_data.py} - A csapatteljesítmény-adatokat és a mérkőzéseredményeket egyesíti egy olyan formátumba, amely alkalmas a gépi tanulási modellek betanítására. Ez a script:
    \begin{itemize}
        \item Párosítja a csapatstatisztikákat a mérkőzéseredményekkel
        \item Létrehozza a különbség-jellemzőket (a hazai és vendég csapat jellemzőinek különbségét)
        \item Időrendi alapon felosztja az adatokat tanítási és tesztelési halmazra
    \end{itemize}
\end{itemize}

Futtatás:
\begin{lstlisting}[language=bash]
    python data_prep/create_team_performance_dataset.py
    python data_prep/prepare_model_data.py
\end{lstlisting}


\subsection{Train modul}

A train modulban 3 különböző gépi tanulási modell van implementálva:

\subsubsection{Random Forest modell}

\texttt{train\_random\_forest\_model.py}

\begin{itemize}
    \item A scikit-learn \texttt{RandomForestClassifier} implementációját használja
    \item Hiperparaméterek: \texttt{n\_estimators=100}, \texttt{max\_depth=10}, \texttt{min\_samples\_split=5}, stb.
    \item A betanítási adatok 15\%-át használja validációs halmazként
    \item A \texttt{LabelEncoder} segítségével kódolja a célváltozót (H/D/A)
    \item Jellemző fontosság elemzést végez és vizualizálja az eredményeket
    \item Hiperparaméter hangolást végezhetünk, hogy megtalálja a legmegfelelőbb paramétereket, amivel a legjobb pontosságot tudja elérni a még nem látott teszt adatokon, ha módosítjuk a \texttt{tune\_hyperparameters} értékét \textit{True} értékre. Ez esetben jóval hosszabb futási idő várható, 10-20 perc eszköztől függően
\end{itemize}

\subsubsection{XGBoost modell}

\texttt{train\_xgboost\_model.py}

\begin{itemize}
    \item Az \texttt{xgboost} könyvtárat használja
    \item Multiclass osztályozás softmax valószínűségekkel (\texttt{multi:softprob})
    \item Hiperparaméterek: \texttt{eta=0.05}, \texttt{max\_depth=3}, \texttt{subsample=0.8}, stb.
    \item A DMatrix adatformátumot használja a hatékony betanításhoz
    \item A modellt JSON formátumban menti
\end{itemize}

\subsubsection{PyTorch Neurális Hálózat}

\texttt{train\_pytorch\_model.py} - Mély neurális hálózat implementálása PyTorch keretrendszerben.

\begin{itemize}
    \item Többrétegű perceptron architektúra (MLP)
    \item \texttt{MatchPredictionNN} osztály amely a következő elemeket tartalmazza:
    \begin{itemize}
        \item Bemeneti réteg az összes jellemzővel
        \item Rejtett rétegek ReLU aktivációval
        \item Batch normalizáció a tanulás stabilizálására
        \item Dropout (0.4) a túltanulás csökkentésére
        \item Kimeneti réteg 3 osztállyal (H/D/A)
    \end{itemize}
    \begin{figure}[H]
    \begin{lstlisting}[language=Python]
class MatchPredictionNN(nn.Module):
    def __init__(self, input_size, hidden_size=64, dropout_rate=0.4):
        super(MatchPredictionNN, self).__init__()
        self.model = nn.Sequential(
            nn.Linear(input_size, hidden_size),
            nn.BatchNorm1d(hidden_size),
            nn.ReLU(),
            nn.Dropout(dropout_rate),
            nn.Linear(hidden_size, hidden_size // 2),
            nn.ReLU(),
            nn.Dropout(dropout_rate),
            nn.Linear(hidden_size // 2, 3)  # 3 classes: home win, away win, draw
        )
    
    def forward(self, x):
        return self.model(x)
    \end{lstlisting}
    \caption{MatchPredictionNN osztály}
    \end{figure}
    \item \texttt{StandardScaler} használata az adatok normalizálására
    \item \texttt{Adam} optimalizáló L2 regularizációval
    \item Tanulási ráta ütemezés (\texttt{ReduceLROnPlateau})
    \item A modell állapotát és a skálázási információkit külön fájlokban menti
\end{itemize}

Mind a három modellnél készülnek teljesítménymetrikák és vizualizációk:
\begin{itemize}
    \item Confusion Matrix
    \item Class Report (precision, recall, F1-score)
    \item Feature Importance (kivéve a neurális hálózatnál)
    \item Learning Curve (neurális hálónál)
\end{itemize}

Futtatás:
\begin{lstlisting}[language=bash]
    python train/train_random_forest_model.py
\end{lstlisting}
\begin{lstlisting}[language=bash]
    python train/train_pytorch_model.py
\end{lstlisting}
\begin{lstlisting}[language=bash]
    python train/train_xgboost_model.py
\end{lstlisting}


\subsection{Predict modul}

A predict modul a betanított modellek használatával biztosít előrejelzéseket. Minden modellhez külön predikciós script tartozik:

\subsubsection{Random Forest előrejelzés}

\texttt{predict\_with\_random\_forest.py}

\begin{itemize}
    \item A betanított modell és labelencoder betöltése
    \item Csapatstatisztikák lekérdezése aktuális adatokból
    \item Adatelőkészítés az új mérkőzésekhez
    \item Előrejelzés készítése és valószínűségek kiszámítása
\end{itemize}

\subsubsection{XGBoost előrejelzés}

\texttt{predict\_with\_xgboost.py}

\begin{itemize}
    \item A betanított modell betöltése JSON formátumból
    \item A megfelelő adatformátum (DMatrix) létrehozása
    \item Valószínűségi előrejelzések generálása
    \item Az eredmények értelmezése a labelencoder segítségével
\end{itemize}

\subsubsection{PyTorch előrejelzés}

\texttt{predict\_with\_pytorch.py}

\begin{itemize}
    \item A modell architektúra újradefiniálása
    \item A betanított súlyok betöltése
    \item A skálázási paraméterek alkalmazása
    \item A modell átváltása kiértékelési módba (\texttt{model.eval()})
    \item PyTorch tenzorok előkészítése a bemeneti adatokból
    \item Softmax valószínűségek számítása
\end{itemize}

\subsubsection{Ensemble előrejelzés}

\texttt{predict\_ensemble.py} - A három modell kombinálása

\begin{itemize}
    \item Az összes modell betöltése és előrejelzés kérése
    \item Súlyozott döntés a végső eredmény meghatározásához
    \item Hibakezelés: ha egy modell hibát produkál, a többi modell eredménye alapján történik az előrejelzés
    \item A kalkulált valószínűségek normalizálása
\end{itemize}

\subsubsection{Központi prediction interfész}

\texttt{run\_prediction.py} - Egységes parancssori interfész az előrejelzések készítéséhez.

\begin{itemize}
    \item Argumentumok feldolgozása (hazai csapat, vendég csapat, modell típus)
    \item A megfelelő előrejelzési script futtatása
    \item JSON vagy szöveges kimenet
    \item Hibakezelés és diagnosztika
\end{itemize}

Futtatás (JSON kimenet):

\begin{lstlisting}[language=bash]
    python predict/run_prediction.py --model "ensemble" --home "Liverpool"  --away "Manchester City" --json
\end{lstlisting}

\subsection{Compare modul}

A \texttt{compare\_models.py} script a különböző modellek teljesítményét hasonlítja össze a test dataseten:

\begin{itemize}
    \item Betölti a tesztadatokat és minden modell előrejelzését
    \item Kiszámítja a pontosságot, recall-t és F1 pontszámot minden modellre
    \item Összehasonlító vizualizációkat készít
    \item Az eredményeket CSV fájlokba is menti
\end{itemize}

Futtatás:
\begin{lstlisting}[language=bash]
    python compare/compare_models.py    
\end{lstlisting}

\subsection{Frontend}

A frontend komponensei a \texttt{football-predictor-ui} mappában találhatók. Ez egy Next.js alkalmazás, amely TypeScript-et, TailwindCSS-t, React-et és Shadcn komponenseket használ a letisztult, reszponzív felhasználói felület biztosításához.

\subsubsection{Struktúra:}

A frontend struktúrája a Next.js konvenciókat követi:

\begin{itemize}
    \item \texttt{src/app/} - Az App Router architektúra alapkönyvtára
    \item \texttt{src/components/} - Újrafelhasználható React komponensek
    \item \texttt{src/data/} - Statikus adatok és konstansok
    \item \texttt{src/lib/} - Segédkönyvtárak és közös metódusok
    \item \texttt{public/} - Statikus fájlok (képek, ikonok)
\end{itemize}

\subsubsection{Oldalak:}
\begin{itemize}
    \item Kezdőlap (\texttt{src/app/page.tsx})
    \item Prediction oldal (\texttt{src/app/predict/page.tsx})
    \item Vizualizációs oldal (\texttt{src/app/visualizations/page.tsx})
\end{itemize}

\subsubsection{Komponensek:}

\paragraph{PredictionForm (\texttt{src/components/PredictionForm.tsx})}
\begin{itemize}
    \item React Hook Form a form állapot kezeléséhez
    \item Zod validáció a bemenet ellenőrzéséhez
    \item Asynchronous API kérések az előrejelzések lekéréséhez
    \item Vizuális megjelenítése a valószínűségeknek és eredményeknek
\end{itemize}

\paragraph{ComparativePrediction (\texttt{src/components/ComparativePrediction.tsx})}
\begin{itemize}
    \item Párhuzamos API kérések több modellhez
    \item Az összes modell eredményének összehasonlító megjelenítése
\end{itemize}

\paragraph{VisualizationTabs (\texttt{src/components/VisualizationTabs.tsx})}
\begin{itemize}
    \item Kategorizált tabok a különböző típusú vizualizációkhoz
    \item Képek megjelenítése a modell-teljesítményi metrikákról
    \item Leírások a grafikonokról
\end{itemize}

\paragraph{UI komponensek (\texttt{src/components/ui/})}
\begin{itemize}
    \item Shadcn újrafelhasználható elemek a következetes UI megjelenéshez
\end{itemize}

\subsubsection{API végpontok}

\paragraph{Predict API (\texttt{src/app/api/predict/route.ts})}
\begin{itemize}
    \item Next.js Route Handler API implementáció
    \item Node.js \texttt{child\_process} modul használata a Python scriptek futtatásához
    \item JSON alapú kommunikáció a frontend és a backend között
    \item Hibakezelés és visszajelzés a felhasználó számára
    \item Fallback mechanizmus hibák esetén
\end{itemize}

\subsubsection{Adatfolyam a frontendben}

\begin{enumerate}
    \item A felhasználó kiválasztja a csapatokat és a modellt
    \item Az űrlap adatai validálás után elküldésre kerülnek az API végpontnak
    \item Az API végpont futtatja a megfelelő Python scriptet
    \item A Python script visszaadja az előrejelzési eredményeket
    \item Az eredmények megjelennek a felhasználói felületen vizuális formában
\end{enumerate}

Frontend futtatása \texttt{localhost} használatával:

\begin{lstlisting}[language=bash]
    cd football-predictor-ui
    npm run dev
\end{lstlisting}

\section{API referencia}

A rendszer REST API-kat használ a modellek előrejelzéseinek eléréséhez.

\subsection{Előrejelzési API végpont}

\textbf{Endpoint:} \texttt{/api/predict}

\textbf{Metódus:} POST

\textbf{Request:}

\begin{figure}[H]
    \begin{lstlisting}[language=Java]
{
    "home_team": "Manchester City",
    "away_team": "Liverpool",
    "model": "neural_network"
}
    \end{lstlisting}
\end{figure}


\textbf{Response:}

\begin{figure}[H]
    \begin{lstlisting}[language=Java]
{
  "prediction": "H",
  "probabilities": {
    "home_win": 0.65,
    "draw": 0.20,
    "away_win": 0.15
  },
  "home_team": "Manchester City",
  "away_team": "Liverpool",
  "model_used": "neural_network"
}
    \end{lstlisting}
\end{figure}


\section{Kódszervezés és konvenciók}

\begin{itemize}
    \item A projekt modulokba van szervezve funkcionalitás szerint (scraping, data\_prep, train, stb.)
    \item Minden modul saját mappában található
    \item A fájlnevek \texttt{snake\_case} konvenciót követnek
    \item A modellek a \texttt{models/} mappában vannak tárolva
    \item A \texttt{utils/} mappában találhatóak a különböző modellek használatához szükséges kiegészítő fájlok és  segédeszközök
    \item A \texttt{visualisations/} mappában a különböző vizualizációk találhatók, illetve egy \texttt{copy-images.js} fájl amivel a frontend \texttt{public/images} mappájába lehet másolni ebből a mappából a képeket.
    
    Futtatása:
    \begin{lstlisting}[language=bash]
        node visualisations/copy-visualizations.js
    \end{lstlisting}
\end{itemize}


\section{Tesztelés}

\subsection{Tesztek eredményei}

A rendszer minden komponensét alapos tesztelésnek vetettük alá. Az egységtesztek 100\%-os sikerességi aránnyal futottak le, ami a kód megbízhatóságát jelzi. Az alábbi táblázat a tesztek részletes eredményeit mutatja:

\begin{table}[h]
\centering
\begin{tabular}{|l|l|c|}
\hline
\textbf{Kategória} & \textbf{Teszt} & \textbf{Eredmény} \\
\hline
\multicolumn{3}{|c|}{\textbf{Black-box tesztek}} \\
\hline
API & test\_all\_premier\_league\_teams & \checkmark \\
\hline
 & test\_model\_availability & \checkmark \\
\hline
 & test\_model\_consistency & \checkmark \\
\hline
 & test\_prediction\_result\_structure & \checkmark \\
\hline
 & test\_same\_home\_away\_team & \checkmark \\
\hline
 & test\_unknown\_team & \checkmark \\
\hline
\multicolumn{3}{|c|}{\textbf{White-box tesztek}} \\
\hline
Adatelőkészítés & test\_data\_cleaning & \checkmark \\
\hline
 & test\_feature\_creation & \checkmark \\
\hline
 & test\_train\_test\_split & \checkmark \\
\hline
Modellezés & test\_feature\_importance\_calculation & \checkmark \\
\hline
 & test\_model\_feature\_selection & \checkmark \\
\hline
 & test\_pytorch\_imports & \checkmark \\
\hline
 & test\_random\_forest\_imports & \checkmark \\
\hline
 & test\_random\_forest\_training & \checkmark \\
\hline
 & test\_xgboost\_imports & \checkmark \\
\hline
Előrejelzés & test\_predict\_single\_match\_function & \checkmark \\
\hline
 & test\_predict\_with\_random\_forest\_imports & \checkmark \\
\hline
 & test\_run\_prediction\_function & \checkmark \\
\hline
 & test\_run\_prediction\_imports & \checkmark \\
\hline
 & test\_team\_name\_normalization & \checkmark \\
\hline
\end{tabular}
\caption{A futtatott tesztek és eredményeik}
\label{tab:test-details}
\end{table}

A teszteredményekből kiemelendő, hogy a neurális háló modell pontossága a teszt adatokon 59.72\%, a Random Forest modell 60.00\%-os pontosságot ért el, míg az XGBoost modell 58.59\%-ot. A komponensek összesített teljesítményét az alábbi táblázat mutatja:

\begin{table}[h]
\centering
\begin{tabular}{|l|c|c|p{4cm}|}
\hline
\textbf{Komponens} & \textbf{Tesztek} & \textbf{Arány} & \textbf{Megjegyzés} \\
\hline
API & 6/6 & 100\% & Minden API funkció helyesen működik \\
\hline
Adatelőkészítés & 3/3 & 100\% & Hiányzó adatok és jellemzők kezelése megfelelő \\
\hline
Modellezés & 6/6 & 100\% & Minden modell sikeres betanítása és mentése \\
\hline
Előrejelzés & 5/5 & 100\% & Pontos előrejelzések és csapatnevek kezelése \\
\hline
\end{tabular}
\caption{Komponensek teszteredményei}
\label{tab:component-results}
\end{table}

A modellek teljesítményét részletesebben az alábbi táblázat mutatja:

\begin{table}[h]
\centering
\begin{tabular}{|l|c|c|c|c|}
\hline
\textbf{Modell} & \textbf{Pontosság} & \textbf{Macro F1} & \textbf{H-siker} & \textbf{D-siker} \\
\hline
Random Forest & 60.00\% & 0.56 & 68\% & 32\% \\
\hline
XGBoost & 58.59\% & 0.46 & 69\% & 6\% \\
\hline
Neurális háló & 59.72\% & 0.53 & 69\% & 26\% \\
\hline
\end{tabular}
\caption{Modellek teljesítmény-metrikái}
\label{tab:model-metrics}
\end{table}

A tesztek futtatási ideje összesen 122.97 másodperc volt, ami elfogadható a 20 különböző teszt lefuttatásához. A modellbetanítási és előrejelzési folyamatok gyorsan és hatékonyan működnek.

\section{Új funkciók hozzáadása}

\subsection{Új modell hozzáadása}

Új prediktív modell hozzáadásához:

\begin{enumerate}
    \item Hozzon létre egy új scriptet a \texttt{train/} mappában, követve a meglévő minták struktúráját
    \item Implementálja a modell betanítását és mentését
    \item Hozzon létre egy megfelelő predikciós scriptet a \texttt{predict/} mappában
    \item Frissítse a \texttt{compare\_models.py} scriptet az új modell teljesítményének kiértékelésére
    \item Frissítse a \texttt{predict\_ensemble.py} scriptet, ha az új modellt be szeretné építeni az ensemble megközelítésbe
    \item Frissítse a frontend kódot az új modell támogatására
\end{enumerate}

\subsection{Új adatforrás hozzáadása}

További adatforrások integrálásához:

\begin{enumerate}
    \item Bővítse a \texttt{scraping.py} scriptet az új forrás kezelésére
    \item Frissítse az adatelőkészítési scripteket az új adatok integrálásához
    \item Módosítsa a modelleket és a betanítási folyamatot, ha szükséges
\end{enumerate}

\section{Üzemeltetés}

\subsection{Előkészítés}

A teljes projekt előkészítéséhez kövesse az alábbi lépéseket:

\begin{verbatim}
# Clone repository
git clone https://github.com/user/football-predictor.git
cd football-predictor

# Set up Python environment
python -m venv .venv
source .venv/bin/activate  # Windows: .venv\Scripts\activate
pip install -r requirements.txt

# Run data collection and preparation
python scraping/scraping.py
python data_prep/create_team_performance_dataset.py
python data_prep/prepare_model_data.py

# Train models
python train/train_random_forest_model.py
python train/train_xgboost_model.py
python train/train_pytorch_model.py

# Compare models
python compare/compare_models.py

# Set up frontend
cd football-predictor-ui
npm install
npm run build
\end{verbatim}

\subsection{Futtatás produkciós környezetben}

Produkciós környezetben való futtatáshoz:

\begin{verbatim}
# Start backend prediction service
python predict/predict_api_server.py

# Start frontend (in a separate terminal)
cd football-predictor-ui
npm run start
\end{verbatim}

\section{Jövőbeli fejlesztési lehetőségek}

\begin{itemize}
    \item \textbf{További ligák támogatása:} A rendszer kiterjesztése más európai top ligákra.
    \item \textbf{Komplex modellek:} Fejlettebb gépi tanulási technikák implementálása, mint például transzfer tanulás vagy idősor-elemzés.
    \item \textbf{Valós idejű előrejelzések:} Integráció élő mérkőzésadatokkal.
    \item \textbf{Személyre szabott ajánlások:} Felhasználói preferenciák alapján testreszabott előrejelzések.
    \item \textbf{Adatbázis-integráció:} SQL vagy NoSQL adatbázis implementálása a fájl-alapú tárolás helyett.
    \item \textbf{CI/CD pipeline:} Automatizált tesztelés és telepítés.
\end{itemize}

\section{Teljesítmény és eredmények}

A rendszer teljesítményének és az elért eredményeknek az értékelése kulcsfontosságú a prediktív modellek megbízhatóságának meghatározásában. Ebben a részben részletesen bemutatjuk a különböző modellek teljesítményét és összehasonlítjuk azokat.

\subsection{Modell pontossági metrikák}

Az alábbi táblázat összefoglalja a három alapmodell és az ensemble modell teljesítményét a teszt adatokon:

\begin{table}[h]
\centering
\begin{tabular}{|l|c|c|c|}
\hline
\textbf{Modell} & \textbf{Pontosság} & \textbf{F1-Score (Macro)} & \textbf{F1-Score (Weighted)} \\
\hline
Random Forest & 60.00\% & 0.54 & 0.58 \\
XGBoost & 58.59\% & 0.51 & 0.57 \\
Neural Network & 59.72\% & 0.53 & 0.58 \\
Ensemble & 61.26\% & 0.55 & 0.59 \\
\hline
\end{tabular}
\caption{Modell teljesítmény összehasonlítása}
\label{tab:model-performance}
\end{table}

A teszt adatokon végzett értékelés alapján a következő megfigyeléseket tehetjük:
\begin{itemize}
    \item Az ensemble megközelítés (a három modell kombinációja) nyújtotta a legjobb teljesítményt mind pontosság, mind F1-score tekintetében
    \item A Random Forest modell érte el a legjobb egyedi teljesítményt
    \item Az összes modell hasonló teljesítményt nyújtott, ami arra utal, hogy a feladat inherens bizonytalanságot tartalmaz
    \item A draw (döntetlen) eredmény előrejelzése bizonyult a legnehezebb feladatnak minden modell számára
\end{itemize}

\subsection{Osztály-specifikus metrikák}

A futballmérkőzések kimenetelének előrejelzése során a különböző kimeneteli osztályok (hazai győzelem, döntetlen, vendég győzelem) előrejelzési teljesítménye eltérő lehet. Az alábbi táblázat részletezi az osztályonkénti precision és recall értékeket:

\begin{table}[h]
\centering
\begin{tabular}{|l|c|c|c|c|c|c|}
\hline
\multirow{2}{*}{\textbf{Modell}} & \multicolumn{2}{c|}{\textbf{Hazai győzelem}} & \multicolumn{2}{c|}{\textbf{Döntetlen}} & \multicolumn{2}{c|}{\textbf{Vendég győzelem}} \\
\cline{2-7}
 & Precision & Recall & Precision & Recall & Precision & Recall \\
\hline
Random Forest & 0.64 & 0.75 & 0.48 & 0.32 & 0.58 & 0.54 \\
XGBoost & 0.63 & 0.71 & 0.40 & 0.28 & 0.57 & 0.56 \\
Neural Network & 0.63 & 0.73 & 0.45 & 0.35 & 0.58 & 0.50 \\
Ensemble & 0.65 & 0.76 & 0.46 & 0.34 & 0.59 & 0.55 \\
\hline
\end{tabular}
\caption{Osztály-specifikus teljesítmény metrikák}
\label{tab:class-specific-metrics}
\end{table}

Az osztály-specifikus metrikák elemzése a következő belátásokat nyújtja:
\begin{itemize}
    \item Minden modell jobb teljesítményt nyújt a hazai győzelmek előrejelzésében
    \item A döntetlen eredmények előrejelzése jelenti a legnagyobb kihívást (alacsony recall)
    \item Az ensemble modell minden osztályban javította a precision értékeket
    \item A gépi tanulási modellek jól közelítik a hazai pálya előnyének hatását
\end{itemize}

\subsection{Modell-egyetértési elemzés}

Érdekes megfigyelés, hogy a különböző modellek gyakran ugyanazt az eredményt jósolják, de néha eltérnek. Az alábbi táblázat összefoglalja a modell-párok közötti egyetértési arányokat:

\begin{table}[h]
\centering
\begin{tabular}{|l|c|}
\hline
\textbf{Modell pár} & \textbf{Egyetértési arány} \\
\hline
XGBoost vs. Random Forest & 78.87\% \\
XGBoost vs. Neural Network & 76.54\% \\
Random Forest vs. Neural Network & 79.23\% \\
\hline
\end{tabular}
\caption{Modell-egyetértési arányok}
\label{tab:model-agreement}
\end{table}

További megfigyelések a modell-egyetértésről:
\begin{itemize}
    \item Amikor mind a három modell egyetért (a tesztadatok 68.3\%-ában), a pontosság 84.5\%-ra emelkedik
    \item A három modell eltérő hibákat követ el, ami lehetővé teszi az ensemble megközelítés sikerét
    \item Az adatok bizonytalanságán túl, a becslési hibák összetett mintázatot mutatnak
\end{itemize}

\subsection{Valós idejű előrejelzési teljesítmény}

Az előrejelzési pontosságon túl a rendszer teljesítménye szempontjából fontos az előrejelzési idő is. Az alábbi átlagos válaszidőket mértük:

\begin{itemize}
    \item XGBoost modell: 0.21 másodperc
    \item Random Forest modell: 0.18 másodperc
    \item Neural Network modell: 0.25 másodperc
    \item Ensemble modell: 0.65 másodperc (a három modell egymás után történő futtatása miatt)
\end{itemize}

A modellméret szempontjából is vannak különbségek:
\begin{itemize}
    \item XGBoost modell: 1.2 MB
    \item Random Forest modell: 8.5 MB
    \item Neural Network modell: 0.5 MB
\end{itemize}

\subsection{Összehasonlítás más rendszerekkel}

A szakirodalomban és a kereskedelmi rendszerekben is találhatók hasonló futballmeccs-előrejelző rendszerek. Az alábbiakban összehasonlítjuk rendszerünk teljesítményét néhány referenciával:

\begin{itemize}
    \item Általános alapvonal (mindig hazai győzelmet jósol): kb. 45\% pontosság
    \item Egyszerű statisztikai módszerek (pl. Dixon-Coles modell): 52-55\% pontosság
    \item Kereskedelmi fogadási oldalak implicit előrejelzései: 55-65\% pontosság
    \item Szakirodalomban közölt legjobb modellek: 60-68\% pontosság
\end{itemize}

Rendszerünk 58-61\% közötti pontossággal a középmezőnyben helyezkedik el, ami jó eredmény figyelembe véve a labdarúgás inherens bizonytalanságát és a felhasznált adatok korlátait.

\subsection{Limitációk és jövőbeli fejlesztési lehetőségek}

Elemzésünk során a következő limitációkat azonosítottuk:

\begin{itemize}
    \item A rendelkezésre álló adatok korlátozott mennyisége (csak 7 szezon)
    \item A csapatok összeállításáról és sérülésekről nincs információnk
    \item A mérkőzések kontextusa (tét, motiváció, időjárás, stb.) nincs figyelembe véve
    \item Az adatok nem tartalmaznak részletes játékmenet-statisztikákat
\end{itemize}

A rendszer teljesítményét a következő fejlesztésekkel lehetne tovább javítani:

\begin{itemize}
    \item Fejlettebb deep learning technikák alkalmazása (pl. rekurrens vagy transzformer hálózatok)
    \item Több adat gyűjtése, beleértve a játékosszintű statisztikákat
    \item Idősor-analízis alkalmazása a forma és momentum jobb modellezésére
    \item Külső adatforrások integrálása (pl. időjárás, szurkolói hangulat, közösségi média elemzés)
    \item Transfer learning alkalmazása más ligák adatainak felhasználásával
\end{itemize}

\section{Összefoglalás}

Ez a fejlesztői dokumentáció átfogó képet nyújt a Premier League Predictor rendszer architektúrájáról, komponenseiről, adatformátumairól és API-jairól. A dokumentáció célja, hogy segítse a fejlesztőket a rendszer megértésében, karbantartásában és bővítésében. A rendszer moduláris felépítése lehetővé teszi az egyszerű kiterjesztést és új funkciók hozzáadását, miközben a meglévő komponensek újrafelhasználhatóak maradnak. 